\subsubsection{Subproceso 2: Selección y Registro de Cursos}



\begin{itemize}
    \item \textbf{Regla del límite de créditos por promedio ponderado (RR-0570-2022):} La cantidad máxima de créditos que el estudiante puede registrar en un semestre se determina según el promedio ponderado de sus dos últimos periodos académicos, conforme a la siguiente escala:
    \begin{itemize}
        \item Promedio $\in [14{,}20;\; 20{,}00]$: hasta 28 créditos.
        \item Promedio $\in [12{,}14;\; 14{,}20)$: hasta 26 créditos.
        \item Promedio $\in [10{,}00;\; 12{,}14)$: hasta 24 créditos.
        \item Promedio $\in [\phantom{0}8{,}00;\; 10{,}00)$: hasta 18 créditos.
        \item Promedio $\in [\phantom{0}0{,}00;\; \phantom{0}8{,}00)$: hasta 15 créditos.
    \end{itemize}

    \item \textbf{Regla de creditaje mínimo obligatorio:} La matrícula regular exige la inscripción de un mínimo de doce (12) créditos semestrales, salvo que al estudiante le resten menos créditos para completar el plan de estudios de su carrera.

    \item \textbf{Regla de prioridad de arrastre:} Es obligatoria la inscripción preferente en las asignaturas desaprobadas del ciclo académico más atrasado y en los cursos clasificados bajo riesgo académico (Grupo Cero), antes de completar la carga con otras asignaturas.

    \item \textbf{Regla de prerrequisito obligatorio:} Queda expresamente prohibida la matrícula en cualquier asignatura cuyos prerrequisitos académicos no figuren como aprobados en el historial del estudiante registrado en el SIGA-ORCE.

    \item \textbf{Regla de incompatibilidad horaria:} Se prohíbe de manera absoluta el cruce de horario en sesiones de laboratorio o prácticas entre dos o más asignaturas. De forma excepcional, el sistema admite hasta un máximo de cuatro (04) horas pedagógicas de cruce exclusivamente entre sesiones teóricas de dos asignaturas distintas.

    \item \textbf{Regla de control de vacantes por concurrencia:} La asignación efectiva de una vacante en una sección queda supeditada a la disponibilidad de cupo existente en el instante exacto en que el estudiante confirma su selección en el portal. El sistema procesa las solicitudes por orden estricto de llegada.
\end{itemize}



